\section{A comparator target and an exact affine minimum}
\label{sec:selector}

Fix $k\ge3$, $d=2k+1$, and $M=2^k-1$. Write an input as
$\xi=(x,y,z)$ with $x,y\in\{-1,1\}^k$ and $z\in\{-1,1\}$, and define
\[
 X=\sum_{i=0}^{k-1}2^i\frac{x_i+1}{2},\qquad
 Y=\sum_{i=0}^{k-1}2^i\frac{y_i+1}{2}.
\]
Our target is
\begin{equation}\label{eq:target}
 f(x,y,z)=+1\quad\Longleftrightarrow\quad z=-1\text{ and }X\ge Y.
\end{equation}
A bitwise comparator followed by a selector gate gives an $O(k)$-size
circuit. Every input bit is essential, so bounded fan-in also gives an
$\Omega(k)$ lower bound on circuit size.

The calibration sample consists of the following inputs, all with $z=-1$.
\begin{enumerate}
\item Positive diagonal points $x=y$: take the all-plus vector, its $k$
one-coordinate flips, and the negatives of these vectors.
\item For each $i$, take two negative carry points with
$x_i=-1,y_i=+1$ and $x_j=+1,y_j=-1$ for $j<i$. Their upper coordinates
$x_j=y_j$, $j>i$, are all plus in one point and all minus in the other.
\end{enumerate}
Discard duplicate points, leaving at most $4k+2$ examples.
For $0\le t\le M$, let $q_t$ have $(X,Y,z)=(t,0,+1)$; all these labels
are negative. Present the calibration, then
$q_k,q_{k+1},\ldots,q_M$, and finally any remaining cube points in a
fixed order. This is a permutation of the cube after duplicates in the
calibration have been removed. The designated prediction rounds will be
$q_{k+1},\ldots,q_M$, of which there are
\begin{equation}\label{eq:mistakes}
 T_k=M-k=2^k-k-1\ge2^{k-1}.
\end{equation}

\begin{lemma}[The affine optimum and a quantitative next-query gap]
\label{lem:selector}
After the calibration and $q_k,\ldots,q_t$, with $k\le t\le M$, the
unique minimum-norm affine vector with signed training margin at least
one is
\begin{equation}\label{eq:At}
 A_t(x,y,z)=-t-(t+1)z+\sum_{i=0}^{k-1}2^i(x_i-y_i).
\end{equation}
Its coefficient norm satisfies
\begin{equation}\label{eq:Qt}
 Q_t=\norm{A_t}_2^2=t^2+(t+1)^2+2\sum_{i=0}^{k-1}4^i,
 \qquad Q_M=\frac83\,4^k-2^{k+1}+\frac13.
\end{equation}
It has absolute margin at least one on the full cube. If $t<M$, then
$A_t(q_{t+1})=1$. Every affine vector $B$ satisfying the same training
constraints and $B(q_{t+1})\le-1$ has
\begin{equation}\label{eq:affine-gap}
 \norm B_2^2\ge Q_t+\frac4{d+1}.
\end{equation}
\end{lemma}
\begin{proof}
Write a competing score as $b+cz+\sum_i a_i x_i+\sum_i h_i y_i$, and put
\[
 u=b-c,\quad v=b+c,\quad w_i=(a_i-h_i)/2,\quad
 s_i=a_i+h_i,\quad U=u-1.
\]
Its squared coefficient norm is
\[
 \frac{u^2+v^2}{2}+2\sum_i w_i^2+\frac12\sum_i s_i^2.
\]
The paired diagonal examples imply $U\ge0$,
$|\sum_i s_i|\le U$, and $|\sum_j s_j-2s_i|\le U$.
Hence $|s_i|\le U$ and $\sum_i|s_i|\le kU$.
Averaging the two carry inequalities at bit $i$ gives
\[
 u-2w_i+2\sum_{j<i}w_j\le-1,
 \qquad
 w_i\ge\sum_{j<i}w_j+\frac{u+1}{2}
       \ge2^i\left(1+\frac U2\right).
\]
At $q_t$, the coordinate contribution is
$2\sum_{\mathrm{bit}_i(t)=1}w_i-\sum_{\mathrm{bit}_i(t)=0}s_i$.
The negative margin therefore implies
\[
 v\le-1-2t-tU+kU\le-(2t+1),
\]
because $t\ge k$. Combining the norm terms gives
\[
 \norm A_2^2\ge\frac{1+(2t+1)^2}{2}+2\sum_i4^i=Q_t.
\]
Equality requires $u=1$, $v=-(2t+1)$, $s_i=0$, and $w_i=2^i$,
which is exactly \eqref{eq:At}. Its scores are
$1+2(X-Y)$ for $z=-1$ and $-(2t+1)+2(X-Y)$ for $z=+1$.
They are odd integers, fit the calibration, and give
$A_t(q_j)=2j-2t-1\le-1$ for $j\le t$. This proves attainment and all
claims except \eqref{eq:affine-gap}.

The training-feasible affine vectors form a closed convex set.
Its minimum-norm point satisfies
$\langle A_t,B-A_t\rangle\ge0$ for every feasible $B$.
The extra query constraint gives
$\langle B-A_t,(1,q_{t+1})\rangle\le-2$.
Since $\norm{(1,q_{t+1})}_2=\sqrt{d+1}$, Cauchy--Schwarz implies
$\norm{B-A_t}_2^2\ge4/(d+1)$. Expanding $\norm B_2^2$ proves the gap.
\end{proof}

We will use the following elementary Gaussian principle. It also
clarifies why all our conclusions concern the full posterior.

\begin{lemma}[A strict energy gap controls posterior mass]
\label{lem:gaussian}
Fix a finite input domain and a finite-dimensional continuous network
with a positive-definite quadratic Gaussian energy $\kappa$.
Let $E$ be the event of strict global margin and a fixed training history.
Suppose a closed parameter set $B$ contains the unwanted part of $E$ and
has $\inf_B\kappa>U$, while $\inf_E\kappa\le U$.
Then $\Pp_\sigma(B\mid E)\to0$ as $\sigma\downarrow0$.
\end{lemma}
\begin{proof}
For a closed set, finite-dimensional Gaussian small-noise bounds give
\[
 \limsup_{\sigma\downarrow0}\sigma^2\log\Pp_\sigma(B)
       \le-\tfrac12\inf_B\kappa.
\]
The event $E$ is open by continuity and its strict inequalities.
Choose a point of $E$ with energy below $\inf_B\kappa$, and then a
sufficiently small ball around that point contained in $E$. The
Gaussian lower bound for this ball gives an exponentially vanishing
ratio to the conditioning-event probability. These inequalities also follow directly
by integrating the Gaussian density: a fixed small ball supplies the
lower bound, while a quadratic-energy sublevel decomposition and Gaussian
tails supply the upper bound. The energy is coercive, so a nonempty
closed constraint set attains its infimum. In particular, if all global
minimizers have one sign function, every other weak-margin sign cell has
strictly larger energy. There are finitely many such cells.
\end{proof}

Whenever $p_t\to0$ on the $T_k$ designated rounds, the claimed forms of
error follow at one common positive scale: choose $\sigma$ small enough
that every $p_t<\min\{\eta/T_k,1/3\}$. A union bound then gives
probability at least $1-\eta$ of a sampled mistake on every designated
round. Majority is wrong and the strict $2/3$ criterion fails there too.
The fixed input order does not depend on sampled predictions.
